\documentclass[11pt]{article}
\usepackage[margin=1in]{geometry}
\usepackage{amsmath,amssymb,amsthm,graphicx,booktabs}
\usepackage[hidelinks]{hyperref}
\newtheorem{theorem}{Theorem}
\newcommand{\Hh}{\ell_2(\mathbb N)}
\newcommand{\ak}{a_k}
\setlength{\parindent}{0pt}
\setlength{\parskip}{0.6em}

\title{Contractivity Cannot Be Removed from the Approximation-Number Truncation Theorem}
\author{Run \texttt{fa\_banach\_001}, lane 14}
\date{9 August 2026}

\begin{document}
\maketitle

\textbf{Status.} Candidate counterexample, likely valid, for human review.

\section{Source question}

Deepesh's Theorem~4 states the following. Let $Y$ be a dual space,
$T\in B(X,Y)$, and $P_n\in B(X)$, $Q_n\in B(Y)$ satisfy
$\|P_n\|\,\|Q_n\|\leq 1$. If
$T_n=Q_nTP_n$ converges to $T$ in the weak-star operator topology, then
$a_k(T_n)\to a_k(T)$ for every $k\in\mathbb N$
\cite[Theorem~4]{Deepesh}.

Section~4 asks whether this conclusion still holds when the norm restrictions
on $P_n$ and $Q_n$ are removed
\cite[Section~4, PDF pp.~13--14]{Deepesh}. The exact source passage spans two
pages:

\begin{center}
\includegraphics[width=\textwidth]{figures/open_problem_crop.png}

\smallskip

\includegraphics[width=\textwidth]{figures/open_problem_crop_continuation.png}
\end{center}

The first sentence of the passage is the target here; the subsequent request
about other $s$-numbers is a separate research direction.

\section{Definition}

For $S\in B(X,Y)$, the $k$th approximation number is
\[
 a_k(S)=\inf\{\|S-F\|:F\in B(X,Y),\ \operatorname{rank}F<k\}.
\]

\section{Proof intuition}

Strong convergence only tests each fixed vector, so it cannot see a block
whose location drifts to infinity. We use finite-rank truncations that are
correct on the first $n$ coordinates, double the next $n$ coordinates, and
discard the rest. The doubled block moves outward, so the truncations converge
strongly to the identity. But its dimension grows: for any fixed $k$, every
rank-less-than-$k$ approximant misses a nonzero direction in that block.
Along that direction the truncation still has size $2$, forcing the $k$th
approximation number to remain $2$ rather than converge to the value $1$ for
the identity.

\section{Counterexample}

Let $H=\Hh$ with standard orthonormal basis $(e_j)_{j\geq1}$, and let $R_m$
denote the orthogonal projection onto
$\operatorname{span}\{e_1,\ldots,e_m\}$. For $n\geq1$, set
\[
 P_n=R_{2n},\qquad E_n=R_{2n}-R_n,\qquad
 Q_n=R_n+2E_n,
\]
and take $X=Y=H$ and $T=I_H$.

\begin{theorem}\label{thm:main}
The operators $P_n$ and $Q_n$ have finite rank and both converge strongly to
$I_H$. Moreover, if $T_n=Q_nTP_n$, then $T_n\to T$ strongly (hence in the
weak-star operator topology), but for every fixed $k\in\mathbb N$,
\[
 a_k(T_n)=2\quad(n\geq k),
 \qquad\text{whereas}\qquad a_k(T)=1.
\]
Consequently the norm restriction in Theorem~4 cannot be removed, even for
finite-rank truncation factors that themselves converge strongly to the
identity.
\end{theorem}

\begin{proof}
Both $P_n$ and $Q_n$ have range contained in $R_{2n}H$, hence have finite
rank. Clearly $P_nx\to x$ for every $x=(x_j)\in H$. The three mutually
orthogonal pieces $R_nH$, $E_nH$, and $(I-R_{2n})H$ show that
\[
 (Q_n-I)x=E_nx-(I-R_{2n})x
\]
and therefore
\[
 \|(Q_n-I)x\|^2
 =\sum_{j=n+1}^{2n}|x_j|^2+\sum_{j>2n}|x_j|^2
 =\sum_{j>n}|x_j|^2\longrightarrow0.
\]
Thus $Q_n\to I$ strongly. Since $Q_nR_{2n}=Q_n$,
\[
 T_n=Q_n I_H P_n=Q_n,
\]
so $T_n\to I_H=T$ strongly. Because $H$ is canonically a dual Hilbert space,
strong convergence implies the weak-star operator convergence required by the
source theorem. Notice also that
\[
 \|P_n\|=1,\qquad \|Q_n\|=2,
\]
so the removed product restriction is violated by the fixed factor $2$.

We next compute the approximation numbers directly. The zero operator gives
$a_k(I_H)\leq1$. Conversely, if $F$ has rank less than $k$, then $\ker F$ is
nonzero (indeed infinite dimensional). For a unit vector $x\in\ker F$,
\[
 \|(I_H-F)x\|=1,
\]
so $a_k(I_H)=1$.

Now fix $k$ and take $n\geq k$. Since $Q_n$ acts as $2I$ on the
$n$-dimensional space $E_nH$, its norm is $2$, and hence $a_k(Q_n)\leq2$.
For any $F$ of rank less than $k$, the restriction
$F|_{E_nH}$ has rank less than $k\leq n=\dim(E_nH)$. It therefore has a
nonzero kernel vector. Choosing a unit vector $x\in E_nH$ with $Fx=0$ gives
\[
 \|(Q_n-F)x\|=\|2x\|=2.
\]
Thus $\|Q_n-F\|\geq2$ for every admissible $F$, so $a_k(Q_n)=2$. Since
$T_n=Q_n$, this proves the displayed failure for every fixed $k$.
\end{proof}

\section{Verification and exact scope}

Every assertion is checked algebraically in the proof. In particular, the
example is not exploiting an unbounded sequence: $\sup_n\|P_n\|=1$ and
$\sup_n\|Q_n\|=2$. It also works simultaneously for all fixed $k$, and both
truncation factors are finite rank and strongly convergent to the identity.

The counterexample answers the literal first question in Section~4. It does
not address the separate request for two-sided convergence theorems for other
$s$-numbers. It also does not determine which intermediate hypotheses---for
example, asymptotic contractivity---might recover the conclusion.

\section{Novelty check}

The bounded search on 9 August 2026 covered the run's registry, solution,
attempt, and proof-gap indexes; exact searches for arXiv:2401.08483, its title,
and the sentence about removing the norm restrictions; and close searches for
approximation-number truncation counterexamples without contractivity. We
checked both arXiv v2 and the official 2025 journal version, which still prints
the question as open. No later exact answer or matching counterexample was
found. This supports, but cannot certify, novelty.

\section{Human-review recommendation}

Accept as a likely valid negative answer after checking that ``the norm
restrictions on $P_n$ and $Q_n$ are removed'' is intended literally relative
to Theorem~4. The finite-rank, strong-convergence construction also satisfies
the natural truncation conditions suggested by the source's motivation.

\begin{thebibliography}{9}
\bibitem{Deepesh}
K. P. Deepesh,
\emph{Approximation results on $s$-numbers of operators},
arXiv:2401.08483v2 (2024); published in \emph{Palestine Journal of
Mathematics} \textbf{14}(3) (2025), 308--318,
\url{https://arxiv.org/abs/2401.08483}.
\end{thebibliography}

\end{document}
